
\documentclass[10pt]{article}
\usepackage[utf8]{inputenc}
\usepackage[T1]{fontenc}
\usepackage{amsmath, amssymb, xcolor, geometry, enumitem, multicol, tcolorbox}

\definecolor{mainblue}{RGB}{0, 102, 204}
\geometry{a4paper, margin=1.2cm, top=1cm, bottom=3cm, footskip=1.2cm}

\begin{document}
\enlargethispage{2.5cm}

% Modern Header
\begin{tcolorbox}[colback=mainblue!5, colframe=mainblue, sharp corners, boxrule=0.5pt]
    \small \textbf{Unit:} undefined \hfill \textbf{Name:} \rule{3cm}{0.4pt} \\
    \small \textbf{Topic:} Variables and Expressions / Order of Operations \hfill \textbf{Date:} \rule{2cm}{0.4pt} \\
    \centering \textbf{\large Challenge (Level 0)}
\end{tcolorbox}

\vspace{0.2cm}

% MCQ Section
\begin{multicols}{2}
\begin{enumerate}[label=\textbf{Q\arabic*.}, leftmargin=*, itemsep=6pt, parsep=0pt]

  \item \small What is the variable in the expression $2x + 3$?
  \begin{enumerate}[label=(\Alph*), leftmargin=0.6cm, nosep, itemsep=2pt]
    \item 2
    \item x
    \item 3
    \item 2x
    \item None of the above
  \end{enumerate}

  \item \small Identify the coefficient in the term $4y$
  \begin{enumerate}[label=(\Alph*), leftmargin=0.6cm, nosep, itemsep=2pt]
    \item y
    \item 4
    \item 4y
    \item None of the above
    \item 0
  \end{enumerate}

  \item \small What is the constant term in the expression $x + 2$?
  \begin{enumerate}[label=(\Alph*), leftmargin=0.6cm, nosep, itemsep=2pt]
    \item x
    \item 2
    \item x + 2
    \item None of the above
    \item 0
  \end{enumerate}

  \item \small Which of the following is an example of a simple expression? $2x + 1$
  \begin{enumerate}[label=(\Alph*), leftmargin=0.6cm, nosep, itemsep=2pt]
    \item Equation
    \item Inequality
    \item Expression
    \item None of the above
    \item Formula
  \end{enumerate}

  \item \small What operation is represented by the symbol $+$ in the expression $2 + 3$?
  \begin{enumerate}[label=(\Alph*), leftmargin=0.6cm, nosep, itemsep=2pt]
    \item Subtraction
    \item Multiplication
    \item Division
    \item Addition
    \item None of the above
  \end{enumerate}

  \item \small Identify the operator in the expression $5 - 2$
  \begin{enumerate}[label=(\Alph*), leftmargin=0.6cm, nosep, itemsep=2pt]
    \item 5
    \item 2
    \item -
    \item None of the above
    \item 5 - 2
  \end{enumerate}

  \item \small What is the term $3x$ an example of?
  \begin{enumerate}[label=(\Alph*), leftmargin=0.6cm, nosep, itemsep=2pt]
    \item Constant term
    \item Variable term
    \item Coefficient
    \item Expression
    \item None of the above
  \end{enumerate}

  \item \small In the expression $2x + 1$, what is the role of the number 2?
  \begin{enumerate}[label=(\Alph*), leftmargin=0.6cm, nosep, itemsep=2pt]
    \item Variable
    \item Constant
    \item Coefficient
    \item Operator
    \item None of the above
  \end{enumerate}

\end{enumerate}
\end{multicols}

\vspace{-0.2cm}
\hrule
\vspace{0.2cm}
\textbf{\small Open Ended Questions (FRQ)}
\begin{enumerate}[label=\textbf{Q\arabic*.}, start=9, itemsep=5pt, topsep=0pt]

  \item \small Draw a simple diagram to illustrate the concept of a variable and a constant in an algebraic expression. Use the expression $3x + 2$ as an example. \vspace{2cm}

  \item \small Explain, in your own words, what the term 'expression' means in algebra. Provide an example of a simple expression and explain its components. \vspace{2cm}

\end{enumerate}

\vfill

% Chef's Tips Footer
\begin{tcolorbox}[colback=blue!5, colframe=mainblue!50, title=\small \textbf{Chef's Tips}, fonttitle=\bfseries, bottom=1mm]
\begin{itemize}[noitemsep, topsep=2pt, leftmargin=0.5cm]
  \item \scriptsize Keep the language and concepts simple, focusing on basic definitions and identification.\item \scriptsize Use visual aids and diagrams to help students understand the concepts of variables, constants, and expressions.\item \scriptsize Emphasize the importance of precise terminology and notation in algebra.
\end{itemize}
\end{tcolorbox}

\end{document}
  